\chapter{Penis Swelling and Injury}

\section*{Definition}
Trauma or swelling involving the penis, ranging from minor contusions to serious injuries including fracture (tunica albuginea rupture), amputation, or strangulation.

\section*{What Happens in Your Body}
Penile fracture occurs when an erect penis is subjected to bending force, causing rupture of the corpus cavernosum tunica albuginea with characteristic popping sound, immediate detumescence, swelling, and ecchymosis. Strangulation (hair tourniquet, constricting bands) causes venous outflow obstruction leading to edema and ischemia. Zipper injuries are common superficial trauma.

\section*{Causes}
\begin{itemize}
  \item Penile fracture (tunica albuginea rupture during intercourse)
  \item Zipper entrapment injury (common in children)
  \item Hair tourniquet or constricting band (strangulation)
  \item Blunt trauma (sports injury, straddle injury)
  \item Penile amputation (rare, surgical emergency)
  \item Iatrogenic (post-circumcision, catheterization)
  \item Insect sting or bite causing angioedema
  \item Paraphimosis (retracted foreskin cannot be reduced)
  \item Peyronie disease (palpable plaque with curvature)
  \item Balanitis or cellulitis with significant edema
\end{itemize}

\section*{When to See a Doctor}
See your doctor if:
\begin{itemize}
  \item Symptoms are severe or getting worse over time
  \item Penile trauma with rapid swelling, inability to urinate, severe pain, blood at urethral meatus, or penile fracture (audible pop during erection)
  \item You have other concerning symptoms such as fever, weight loss, or bleeding
\end{itemize}

\section*{Related Symptoms}
\begin{itemize}
  \item Penis pain
  \item Blood in urine
  \item Difficulty urinating
  \item Swelling
  \item Bruising
\end{itemize}
\textbf{Important:} If this symptom persists, your doctor will check for serious underlying causes.

\section*{Self-Care / Home Management}
\begin{itemize}
  \item Rest and avoid activities that make it worse.
  \item Maintain adequate hydration and a balanced diet.
  \item Over-the-counter remedies may provide symptomatic relief (consult a pharmacist or doctor).
  \item Monitor symptoms and seek professional advice if they persist or worsen.
  \item Avoid self-medicating with prescription medications without medical guidance.
\end{itemize}

\section*{How Your Doctor Will Evaluate You}
\begin{itemize}
  \item A thorough history and physical examination are the first steps in evaluation.
  \item Laboratory tests (blood work, urinalysis) may be ordered based on clinical suspicion.
  \item Imaging studies (X-ray, ultrasound, CT, MRI) may be indicated when structural pathology is suspected.
  \item Specialist referral is arranged when the diagnosis remains uncertain or requires specific expertise.
\end{itemize}

\section*{Treatment Options}
\begin{itemize}
  \item Treatment targets the underlying cause identified through diagnostic evaluation.
  \item Supportive care and symptom management are provided while awaiting definitive therapy.
  \item Medications, lifestyle modifications, and physical therapies may be prescribed as appropriate.
  \item Follow-up ensures treatment effectiveness and allows timely adjustments to the care plan.
\end{itemize}

\clearpage